% META: {"id": "TCOMPL-BAYES-CO-007", "chapitre": "TCOMPL-INFERENCE-BAYESIENNE", "type_objet": "corrige", "exercice_ref": "TCOMPL-BAYES-EX-007", "capacites_codes": ["C3"], "sources_inspiration": [], "status": "generated"}
% BEGIN-VERIFY
% from sympy import Rational, N
% pm, sens, faux = Rational(2, 100), Rational(95, 100), Rational(4, 100)
% pt = pm * sens + (1 - pm) * faux
% vpp = pm * sens / pt
% assert sens != vpp
% assert float(N(sens)) > 0.9 and float(N(vpp)) < 0.4
% # Sans la prevalence, la VPP n'est pas determinee : elle varie de 0 a 1.
% for prevalence in (Rational(1, 1000), Rational(1, 2), Rational(9, 10)):
%     p = prevalence * sens + (1 - prevalence) * faux
%     v = prevalence * sens / p
%     assert 0 < float(N(v)) < 1
% assert float(N(Rational(1, 1000) * sens / (Rational(1, 1000) * sens + Rational(999, 1000) * faux))) < 0.03
% END-VERIFY
\begin{corrige}{TCOMPL-BAYES-EX-007}
\textbf{1.} $P_M(T^+)=0{,}95$ est la sensibilité : la probabilité que le test
détecte la maladie \emph{chez un malade}. $P_{T^+}(M)\approx0{,}33$ est la
valeur prédictive positive : la probabilité d'être malade \emph{sachant que le
test est positif}. Les deux conditionnements sont inversés.

\textbf{2.} Le patient confond $P_M(T^+)$ et $P_{T^+}(M)$. Ces deux
probabilités ne sont égales que dans des cas très particuliers.

\textbf{3.} Il faut connaître la \emph{prévalence} $P(M)$ : c'est elle qui,
dans la formule de Bayes, fait chuter la valeur prédictive lorsque la maladie
est rare.
\end{corrige}
